\section{Stokes lifting and differentiated source interactions}\label{r:joined-sec:unrestricted}

The prescribed force determines a linear Stokes response with
finite norms at every fixed order. Separating this response from the
velocity makes the additional nonlinear interactions explicit.
Let $z$ solve
\begin{equation}\label{r:joined-eq:stokeslift}
z_t-\nu\Delta z=\PP f,\qquad z(0)=0,
\qquad z(t)=\int_0^t e^{\nu(t-s)\Delta}\PP f(s)\,ds.
\end{equation}
The heat semigroup is a contraction in every $H^q$. For $Z_n=\partial_t^nz$,
\begin{equation}\label{r:joined-eq:liftjets}
\begin{aligned}
Z_{n+1}&=\nu\Delta Z_n+\PP F_n,\\
Z_n(0)&=\sum_{j=0}^{n-1}(\nu\Delta)^{n-1-j}\PP F_j(0)
\quad(n\ge1),\\
\sup_{t\le T}\norm{Z_n(t)}_{H^q}
&\le\norm{Z_n(0)}_{H^q}+\int_0^T\norm{F_n(s)}_{H^q}\,ds.
\end{aligned}
\end{equation}
Every norm in this estimate depends on finitely many derivatives of
the prescribed force. Set \(v=u-z\). This decomposition separates the
known Stokes response from the remainder of the forced solution.

With $B(a,b)=\PP[(a\cdot\nabla)b]$, the equation for $v$ is
\begin{equation}\label{r:joined-eq:liftequation}
v_t-\nu\Delta v+B(v,v)+B(z,v)+B(v,z)+B(z,z)=0.
\end{equation}
Put $W_n=\partial_t^nv$, so $V_n=W_n+Z_n$. Every temporal row becomes
\begin{equation}\label{r:joined-eq:liftrow}
\begin{aligned}
W_{n+1}-\nu\Delta W_n
+\sum_{a=0}^n\binom na\bigl\{
&B(W_a,W_{n-a})+B(Z_a,W_{n-a})\\
&+B(W_a,Z_{n-a})+B(Z_a,Z_{n-a})\bigr\}=0.
\end{aligned}
\end{equation}
The binomial self-interaction is retained in full. The other three terms contain a prescribed lift factor. If $\nabla\phi=\QQ f$ and $\pi=p-\phi$, pressure is retained by
\begin{equation}\label{r:joined-eq:liftpressure}
-\Delta\partial_t^n\pi
=\sum_{a=0}^n\binom na
\partial_i(W_a+Z_a)_j\,
\partial_j(W_{n-a}+Z_{n-a})_i.
\end{equation}

\begin{theorem}[Finite-order estimate for the Stokes decomposition]\label{r:joined-thm:lift}
For a fixed finite $N$ and real $s\ge0$, define
\begin{equation}\label{r:joined-eq:liftenergy}
\begin{aligned}
\mathscr E&=\tfrac12\sum_{n=0}^N\norm{W_n}_{H^s}^2,
&\mathscr D&=\sum_{n=0}^N\norm{\nabla W_n}_{H^s}^2,\\
\mathscr Z&=\sum_{a=0}^N\binom Na\norm{Z_a}_{H^{s+2}},
&G_n&=\sum_{a=0}^n\binom na B(Z_a,Z_{n-a}),\\
\mathscr G^2&=\sum_{n=0}^N\norm{G_n}_{H^{s-1}}^2.
\end{aligned}
\end{equation}
Let the signed self-production be
\begin{equation}\label{r:joined-eq:selfproducer}
\mathscr P_{N,s}(W)
=-\sum_{n=0}^N\ip{W_n}{
 \sum_{a=0}^n\binom na B(W_a,W_{n-a})}_{H^s}.
\end{equation}
On every strict classical interval,
\begin{equation}\label{r:joined-eq:unrestrictedrowestimate}
\mathscr E'+\frac\nu2\mathscr D
\le\mathscr P_{N,s}(W)
+\left(\nu+\frac{8C_s^2}{\nu}\mathscr Z^2\right)\mathscr E
+\frac{\mathscr G^2}{\nu},
\end{equation}
where $C_s=2^{s/2}\pi(2\pi)^{-3/2}$. The two added coefficients have finite time integrals determined entirely by the prescribed force. No smallness condition is used.
\end{theorem}
\begin{proof}
The Fourier weight inequality
$\langle\xi\rangle^s\le2^{s/2}\langle\eta\rangle^s
\langle\xi-\eta\rangle^s$
and Cauchy--Schwarz give
\begin{equation*}
\norm{ab}_{H^s}\le C_s\norm a_{H^{s+2}}\norm b_{H^s},
\end{equation*}
using $\int\langle\eta\rangle^{-4}d\eta=\pi^2$.
Since both factors are divergence-free, the divergence form of $B$ gives
$\norm{B(a,b)}_{H^{s-1}}\le\norm{a\otimes b}_{H^s}$.
In either cross term, the factor placed at height $s+2$ is $Z$.

Let $L_n$ be the sum of the two cross interactions in \eqref{r:joined-eq:liftrow}. Discrete convolution and $\binom na\le\binom Na$ give
\begin{equation*}
\left(\sum_{n=0}^N\norm{L_n}_{H^{s-1}}^2\right)^{1/2}
\le2C_s\mathscr Z\sqrt{2\mathscr E}.
\end{equation*}
Pairing through heights $s-1$ and $s+1$, all added terms are bounded by
\begin{equation*}
\sqrt{\mathscr D+2\mathscr E}
\left(2C_s\mathscr Z\sqrt{2\mathscr E}+\mathscr G\right).
\end{equation*}
Young's inequality bounds this by
$\frac\nu2\mathscr D+\nu\mathscr E
+8C_s^2\mathscr Z^2\mathscr E/\nu+\mathscr G^2/\nu$.
The remaining pairing is exactly \eqref{r:joined-eq:selfproducer}, proving
\eqref{r:joined-eq:unrestrictedrowestimate}. Finally,
$\mathscr G\le C_s\mathscr Z(\sum_{n\le N}\norm{Z_n}_{H^s}^2)^{1/2}$,
so \eqref{r:joined-eq:liftjets} bounds every added coefficient.
\end{proof}

The added interactions have coefficients determined by the source,
and their contribution is at most linear in \(\mathscr E\).
At \(N=s=0\), incompressibility also cancels the self-production.
At higher orders, the signed nonlinear term in
\eqref{r:joined-eq:selfproducer} remains. The estimate controls the
source-dependent perturbation of the energy equation without giving
an independent bound for that nonlinear term.

The decomposition must also retain the original initial derivatives.
For \(a_n=V_n(0)\), the recurrence \eqref{r:eq:initial} reads
\begin{equation}\label{r:joined-eq:initialjet}
a_{n+1}=\nu\Delta a_n
-\PP\sum_{b=0}^{n}\binom nb(a_b\cdot\nabla)a_{n-b}
+\PP F_n(0).
\end{equation}
Thus the initial derivatives remain fixed by \(u_0\) and the
source derivatives at zero. Full-interval regularity for arbitrary
forcing is supplied by Corollary~\ref{r:thm:supplied-forced}, as a direct
consequence of the global forced theorem proved in Part~II.


\subsection{The force inside differentiated nonlinear production}

The explicit work term does not contain every occurrence of the source after time differentiation. The nonlinear production is evaluated on a forced velocity, so differentiating it inserts the force into each velocity slot. At the second critical-height row this entire insertion can also be controlled from the kinetic budget.

Write $E_s(u)=\frac12\norm u_{\dot H^s}^2$, $g=\PP f$,
$\mathcal A(u)=\nu\Delta u-B(u,u)$ and
$\mathcal P_s(u)=-\langle\Lambda^su,\Lambda^sB(u,u)\rangle$.
The vector field $\mathcal A$ is the unforced right side evaluated at the current velocity; the actual trajectory obeys $u_t=\mathcal A(u)+g$.
Let $\mathcal J_s=D\mathcal P_s(u)[g]$. Differentiating its three velocity slots gives
\begin{equation}\label{r:joined-eq:sourceinsertion}
\begin{aligned}
\mathcal J_s={}&-\langle\Lambda^sg,\Lambda^s((u\cdot\nabla)u)\rangle\\
&-\langle\Lambda^su,\Lambda^s((g\cdot\nabla)u)\rangle\\
&-\langle\Lambda^su,\Lambda^s((u\cdot\nabla)g)\rangle .
\end{aligned}
\end{equation}
The projection has been removed only inside global pairings with solenoidal fields.

Define $U=\sup_{t<T}\norm{u(t)}_2$,
$D=\norm{\nabla u}_{L^2(0,T;L^2)}$, and
\begin{equation*}
a_\rho(t)=(2\pi)^{-3/2}\int_{\R^3}|\xi|^\rho
|\widehat g(t,\xi)|\,d\xi,\qquad
h(t)=\norm{u(t)}_{\dot H^{1/2}}.
\end{equation*}
Kinetic energy bounds $U$ and $D$ from the prescribed data. Every $a_\rho$ used below is bounded by a sufficiently high source Sobolev norm.

\begin{proposition}[Complete critical source insertion]\label{r:joined-prop:sourceinsertion}
Without a source smallness condition,
\begin{equation}\label{r:joined-eq:insertionbound}
|\mathcal J_{1/2}|\le
2a_1h^2+a_{3/2}Uh+a_2U^2.
\end{equation}
In particular, with $a_\rho^*=\sup_{t\le T}a_\rho(t)$,
\begin{equation}\label{r:joined-eq:insertionL1}
\norm{\mathcal J_{1/2}}_{L^1(0,T)}
\le2a_1^*U\sqrt T\,D
+a_{3/2}^*U^{3/2}T^{3/4}D^{1/2}
+a_2^*TU^2.
\end{equation}
\end{proposition}
\begin{proof}
In the first pairing of \eqref{r:joined-eq:sourceinsertion}, move all spatial derivatives onto $g$; its absolute value is bounded by $a_2U^2$.
In the second pairing the undifferentiated transport cancels by $\diver g=0$. Its commutator obeys
\begin{equation*}
\norm{[\Lambda^{1/2},g\cdot\nabla]u}_2\le a_1h,
\end{equation*}
because
$\bigl||\xi|^{1/2}-|\eta|^{1/2}\bigr|\,|\eta|
\le|\xi-\eta|\,|\eta|^{1/2}$
in its Fourier kernel. Young's convolution inequality proves the bound.
For the final pairing, splitting the half-derivative Fourier weight gives
\begin{equation*}
\norm{\Lambda^{1/2}((u\cdot\nabla)g)}_2
\le a_1h+a_{3/2}U.
\end{equation*}
Summing gives \eqref{r:joined-eq:insertionbound}. Finally, Fourier Cauchy--Schwarz gives $h^2\le U\norm{\nabla u}_2$. Integrating this inequality and its square root proves \eqref{r:joined-eq:insertionL1}.
\end{proof}

The direct work terms can also be estimated using only kinetic energy and source norms. Set
$W_s=\langle u,\Lambda^{2s}g\rangle$ and $k_s=\Lambda^{2s}g$.
Moving derivatives onto the prescribed test field yields
\begin{equation}\label{r:joined-eq:workderivative}
\begin{aligned}
W_s'={}&\langle u,(\partial_t+\nu\Delta)k_s\rangle\\
&+\int_{\R^3}u_i u_j\,\partial_j(k_s)_i\,dx
+\langle g,k_s\rangle .
\end{aligned}
\end{equation}
Each term is uniformly bounded from $U$ and smooth source norms. The pressure pairing vanishes because $k_s$ is divergence-free.

Let $H(u)=\frac12\norm u_{\dot H^{1/2}}^2$ and write
$L_{\mathcal A}\Phi=D\Phi[\mathcal A]$ for differentiation of a scalar functional along the unforced vector field. Along the forced history the exact second-height equation is
\begin{equation}\label{r:joined-eq:generatorcomparison}
\begin{aligned}
H''&=L_{\mathcal A}^2H(u)+R_f,\\
R_f&=\mathcal J_{1/2}+W_{1/2}'-2\nu W_{3/2},
\qquad R_f\in L^1(0,T).
\end{aligned}
\end{equation}
Thus the source correction hidden inside nonlinear differentiation is controlled along with the explicit work terms. The remaining functional is
\begin{equation}\label{r:joined-eq:autonomousgenerator}
L_{\mathcal A}^2H(u)
=4\nu^2E_{5/2}(u)
+D\mathcal P_{1/2}(u)[\mathcal A(u)]
-2\nu\mathcal P_{3/2}(u).
\end{equation}
The leading viscous coefficient remains positive.
The remainder depends on the nonlinear evolution of \(u\).
Subtracting the double time integral of \(R_f\) corrects the scalar
energy identity, but the velocity still satisfies
\(u_t=\mathcal A(u)+g\). Thus neither this scalar correction nor the
Stokes decomposition identifies the forced solution with an unforced
one. The estimates establish bounded source contributions at the
stated orders; full-interval control of the nonlinear remainder is
a separate assertion.
